% Study sheet: Volume II, Chapter 17, Conclusion.
% Numbers: \val{...} from conclusion_studysheet_numbers.tex, written by conclusion_studysheet_numbers.py.
\documentclass{studysheet}
\input{conclusion_studysheet_numbers}
\sheetvolume{2}
\sheetchapter{17}{Conclusion}
\sheetsources{Numbers from the mlsysim registry and the chapter's modeled scenarios (GPU-only failures); questions marked assumed use new inputs.}

\begin{document}
% ===================================================================== FRONT: LEARN
\frontside
\begin{bigidea}
At fleet scale a better model or accelerator alone does not make a better service. The volume's thread:
the \hl{C\textsuperscript{3}~taxonomy} names every cost as compute, communication or coordination, and
\hl{the Fleet Law} prices them in step time; six principles locate where that price is paid. The volume
closes on one method: find the binding term, measure its cost, keep displaced costs visible.
\end{bigidea}

\begin{leftcol}
\sheetsec{Key ideas}
\begin{ideas}
\idea{Scale can move the binding constraint off the node.}{It moves into the fabric, recovery, the service
  contract and obligations.}
\idea{Find the binding term; never assume it.}{Barrier idle time is ambiguous: per-rank traces split late
  arrival (coordination) from slow transfer (communication).}
\idea{Overhead is displaced, not removed.}{Gradient compression cuts bytes but moves cost into convergence
  risk; keep evidence that quality still holds.}
\idea{For fixed work, growth buys calendar time.}{A bigger fleet finishes sooner but spends more
  accelerator-hours.}
\idea{Obligations spend the existing terms.}{Verification uses compute, approval gates add coordination,
  and every second of step time costs joules.}
\idea{Judge a change at the fleet boundary.}{Halving link energy per bit at fixed traffic saves at most half
  the links' share of fleet power (Amdahl on energy).}
\end{ideas}

\vspace{10pt}
\sheetsec{The diagnostic procedure}
{\centering
\begin{tikzpicture}[sheet, every node/.style={nbox, text width=58pt, align=center, minimum height=20pt}]
  \node (a) at (38,20) {observed symptom};
  \node (b) at (122,20) {falsifiable metric};
  \node[accbox, text width=58pt, align=center, minimum height=20pt] (c) at (206,20) {binding C\textsuperscript{3} term};
  \node (d) at (206,-18) {stack row that owns the fix};
  \node (e) at (122,-18) {cost the fix displaces};
  \node[hotbox, text width=58pt, align=center, minimum height=20pt] (f) at (38,-18) {evidence to keep};
  \draw[arr] (a) -- (b); \draw[arr] (b) -- (c); \draw[arr] (c) -- (d); \draw[arr] (d) -- (e); \draw[arr] (e) -- (f);
\end{tikzpicture}\par}

\vspace{10pt}
\sheetsec{Six principles and their metrics}
\begin{sheettable}{@{}L L L@{}}
\thead{Principle} & \thead{Key metric} & \thead{Practice it demands}\\ \midrule
Infrastructure determines capability & FLOP/s, bandwidth, power & co-design hardware with algorithms\\
Communication can dominate & exposed comm time & expect bottlenecks to migrate\\
Failure is routine & MTBF, checkpoint overhead & treat failure as steady state\\
Serving is bounded by its contract & p99 within the SLO & reserve headroom\\
Responsibility constrains design & subgroup error, privacy budget & instrument first\\
Sustainability is a first-order cost & energy per query, carbon & compound small gains\\
\end{sheettable}
\end{leftcol}%
\begin{rightcol}
\sheetsec{Formulas you will use}
\fcard{Fleet Law: the price}%
  {\DistributedStepTime(\DeviceCount) = \under{\frac{\IronLawComputeTime}{\DeviceCount}}{compute}
   + \under{\CommunicationTime(\DeviceCount)}{communication} + \under{\SynchronizationTime(\DeviceCount)}{coordination}
   - \under{\OverlappedTime}{hidden}}%
  {On $\DeviceCount$ devices, only the compute term shrinks. Recovery pauses add to $\SynchronizationTime$.}{}{}
\fcard{Recovery cost of a bigger fleet}%
  {\SystemMtbf = \frac{\MeanTimeToFailure}{\DeviceCount}, \qquad
   \CheckpointOverheadFraction \approx \sqrt{\frac{2\,\CheckpointWriteTime}{\SystemMtbf}}}%
  {Per-device MTTF over $\DeviceCount$ devices; at the Young-Daly interval, checkpoint-and-rework overhead
   (write time $\CheckpointWriteTime$) grows as $\sqrt{\DeviceCount}$. With uninterrupted step time
   $\DistributedStepTime$, fixed work takes $\propto \DistributedStepTime / (1-\CheckpointOverheadFraction)$:
   break-even step speedup $= (1-\CheckpointOverheadFraction)$ before $\div$ after. Accelerator-hours
   $= \DeviceCount \times$ run time.}%
  {Independent GPU failures; the job needs every GPU.}{}
\fcard{What a doubling buys}%
  {\ScalingLawLoss \propto \TrainingCompute^{-\ScalingExponent}
   \;\Rightarrow\; \text{doubling multiplies } \ScalingLawLoss \text{ by } 2^{-\ScalingExponent}}%
  {Fitted exponent \val{kaplanalpha}: loss $\times$\val{lossfactor} (\val{losscut}\% less) per doubling;
   same fraction, smaller absolute gain, full price.}%
  {Only inside the fitted range.}{}

\vspace{10pt}
\sheetsec{Two decisions, separate evidence}
\begin{sheettable}{@{}l L@{}}
\thead{Decision} & \thead{Passes when}\\ \midrule
Accept the expansion & measured step speedup beats break-even\\
Release the model & gates set before exposure: p99 under production load with rollback headroom;
  provenance of data, code, weights; bounded disclosure; energy accounted\\
\end{sheettable}
\end{rightcol}

% ===================================================================== BACK: PRACTICE
\backside
\begin{answers}
\ans{Finish}{\val{mttf} $\div$ \val{quadgpus} $=$ \val{quadmtbf} h $=$ \val{quadmtbfsec} s; $\sqrt{\val{youngdaly} \times \val{writes} \div \val{quadmtbfsec}}$
  $=$ \val{quadover}\% (\val{quadovergrowth}$\times$ the reference: $\sqrt{\DeviceCount}$). Break-even \val{refkeep} $\div$ \val{quadkeep} $=$
  \val{quadbreakeven}$\times$; at the ideal \val{quadfactor}$\times$ step, accelerator-hours rise \val{quadhoursup}\%.}
\ans{1}{(a) Aligned per-rank traces: arrival skew versus exposed collective time. (b) Communication (communication
  can dominate); expect the bottleneck to migrate, so measure again after the fix. (c) Convergence and accuracy
  risk; keep the evaluation record showing quality still holds.}
\ans{2}{(a) Yes for time: \val{qspeedup} $>$ \val{breakeven}. Accelerator-hours $\times$ \val{factor} $\div$ \val{qspeedup} $\times$ \val{breakeven}
  $=$ \val{qhoursratio}, \val{qhoursup}\% more for the same work. (b) No: a release gate fails. It returns to the row that owns it (assurance)
  and the model ships only after the revised candidate passes (two-decisions table). (c)
  \val{lossfactor}$^{\val{fourdoublings}}$ $=$ \val{fourlossfactor}, a \val{fourlosscut}\% cut, inside the fitted range only (doubling card).}
\ans{3}{At most \val{qlinkcut}\% of \val{qlinkshare}\% $=$ \val{qfleetsave}\% of fleet power: judge at the fleet boundary (Amdahl applied to
  energy), and only while traffic stays fixed.}
\ans{4}{\val{llamadays} $\times$ \val{hoursperday} h $\div$ \val{llamastops} $\approx$ \val{llamahours} h between stops. Failure is routine: treat it as steady state
  (principles table): budget checkpoint writes and rollback rework (recovery card).}
\end{answers}

\begin{leftcol}
\sheetsec{Worked example}
\begin{practicebox}\raggedright\boxtitle{Does doubling the cluster pay for its recovery cost?}
Archetype A: a \val{params}B model on \val{refgpus} GPUs grows to \val{expgpus}. Per-GPU MTTF \val{mttf} h; the
\val{ckpttb} TB checkpoint takes \val{writes} s to write.\par\vspace{5pt}
{\centering
\begin{tikzpicture}[sheet]
  \pgfmathsetmacro\sx{190/4096} \pgfmathsetmacro\sy{48/4}
  \draw[dim] (30,0) -- (225,0); \draw[dim] (30,0) -- (30,52);
  \foreach \gx/\gl in {1024/{1,024},2048/{2,048},4096/{4,096}}{
    \draw[dim] ({30+\gx*\sx},0) -- ({30+\gx*\sx},-2.5);
    \node[small, anchor=north] at ({30+\gx*\sx},-4) {\gl};}
  \foreach \yv in {0,2,4}{\node[small, anchor=east] at (27,{\yv*\sy}) {\yv\%};}
  \node[small, anchor=north east] at (225,-13) {GPUs};
  \draw[accent, line width=1pt] plot[domain=0:4096, samples=60, smooth] ({30+\x*\sx},{\val{curvecoef}*sqrt(\x)*\sy});
  \filldraw[fill=white, draw=accent] ({30+1024*\sx},{\val{refoverraw}*\sy}) circle (2.2pt);
  \filldraw[fill=accent, draw=accent] ({30+2048*\sx},{\val{expoverraw}*\sy}) circle (2.2pt);
  \node[note, anchor=south east] at ({27+1024*\sx},{\val{refoverraw}*\sy+3}) {\val{refover}\%};
  \node[note, anchor=south east] at ({27+2048*\sx},{\val{expoverraw}*\sy+3}) {\val{expover}\%};
  \node[lab, anchor=west] at (150,14) {overhead grows as $\sqrt{\DeviceCount}$};
\end{tikzpicture}\par}
\begin{steps}
\item Interval: \val{mttf} $\div$ \val{refgpus} $=$ \res{\val{refmtbf} h}; $\div$ \val{expgpus} $=$ \res{\val{expmtbf} h}.
\item Overhead: $\sqrt{\val{youngdaly} \times \val{writes}\text{ s} \div \val{refmtbfsec}\text{ s}}$ $=$ \res{\val{refover}\%}; on
  \val{expgpus}, \res{\val{expover}\%} ($\times$\val{overgrowth} $=\sqrt{\val{factor}}$).
\item Break-even step speedup: \val{refkeep} $\div$ \val{expkeep} $=$ \res{\val{breakeven}$\times$}: recovery alone does not block it.
\item At the ideal \val{factor}$\times$ step, accelerator-hours rise \res{\val{hoursup}\%}.
\end{steps}
\takeaway{Recovery alone barely moves break-even; the measured step speedup, set by exposed communication, decides.}
\end{practicebox}

\vspace{10pt}
\sheetsec{Your turn: finish it}
\begin{practicebox}\raggedright\boxtitle{\assumed{} quadruple instead, to \val{quadgpus} GPUs}
\textbf{Given:} interval \val{mttf} $\div$ \val{quadgpus} $=$ \val{quadmtbf} h.\par
\textbf{You:} find the overhead, the break-even step speedup, and the accelerator-hour rise at the ideal
step.\par\vspace{40pt}
\end{practicebox}
\end{leftcol}%
\begin{rightcol}
\sheetsec{Common mistakes}
\mistake{``A faster model is automatically a better system.''}%
  {In a barrier-dominated step, one worker at \val{slowspeed}\% of peer speed cuts throughput
   \val{slowdrop}\%; faster kernels elsewhere cannot shorten that wait.}{}
\mistake{``Each server meets its p99, so requests are fast.''}%
  {A request touching \val{fanservers} servers has a \val{fanhit}\% chance that at least one answers
   beyond its p99.}{}
\mistake{``We will add the fairness and audit checks at release review.''}%
  {A disparity never measured cannot be detected, attributed, or remedied: instrument first.}{}

\vspace{2pt}
\sheetsec{Check yourself}
\begin{qs}
\item After the expansion, per-step MFU dropped and accelerators idle at gradient barriers. (a) What
  evidence separates slow transfer from late arrival? (b) Traces show little skew, longer exposed transfer:
  which term binds, and which practice applies after the fix? (c) Gradient compression displaces what cost, and what evidence is kept?
\item \assumed{} on the doubled cluster the measured uninterrupted step speedup is \val{qspeedup}$\times$, and the
  release review lacks evidence bounding output disclosure. (a) Should the expansion proceed, and at what
  accelerator-hour cost? (b) Can the model ship? (c) If the team instead quadruples compute, what loss cut
  does the fit illustrate?
\item \assumed{} fabric links draw \val{qlinkshare}\% of fleet power; new optics cut energy per bit
  \val{qlinkcut}\% at fixed traffic. What fleet power saving can you claim, and by which rule?
\item One Llama 3 run on \val{llamagpus} GPUs saw \val{llamastops} unexpected interruptions in
  \val{llamadays} days, yet kept over \val{llamakept}\% of its time as useful training. How often did it stop, which
  principle does that show, and what must the design budget for?
\end{qs}
\end{rightcol}
\end{document}
